\section{Chapter~\ref{sec:sleep}. Sleep}

The sleep modes, the replays, and the shrinking of synapses are measured by neuron recordings, microdialysis, and electron microscopy; the sleep schedule and the slow swings are described by the book's models, while the purposes of sleep are mostly hypotheses.

{\footnotesize
\setlength{\tabcolsep}{3pt}\renewcommand{\arraystretch}{1.15}
\begin{longtable}{>{\raggedright}p{0.5cm}>{\raggedright}p{4.0cm}>{\raggedright}p{5.6cm}>{\raggedright}p{2.3cm}>{\raggedright\arraybackslash}p{1.7cm}}
\toprule
No. & Claim & Experiment and what was measured & Source & Status \\
\midrule\endfirsthead
\toprule
No. & Claim & Experiment and what was measured & Source & Status \\
\midrule\endhead
1 & Cortical neurons are not silent in sleep; the operating mode changes & Long recordings of cortical neurons in rats: in slow-wave sleep the rate stays nonzero, and periods of activity alternate with periods of silence; after long wakefulness the rate is higher in all states, after sleep it is lower & \cite{Vyazovskiy2009} & measured \\
2 & \nt{ACET} is high in the day and in REM sleep, low in slow-wave sleep (Table~\ref{tab:sleepmodes}) & Microdialysis in the cortex of freely moving cats: in wakefulness acetylcholine release is $100$--$175\%$ higher than in slow-wave sleep, and in REM sleep it is about at the level of quiet wakefulness & \cite{Marrosu1995} & measured \\
3 & \nt{NORA} and \nt{SERO} quiet down in slow-wave sleep and are nearly silent in REM sleep & Recordings of single \nt{NORA} neurons in rats and \nt{SERO} neurons in cats across sleep stages: the rate is highest in wakefulness, lower in slow-wave sleep, and nearly zero in REM sleep & \cite{AstonJonesBloom1981,McGintyHarper1976} & measured \\
4 & In REM sleep the muscles are switched off by inhibition through \nt{GABA} and \nt{GLYC} & In rats the activity of the \nt{ACET} neurons of a jaw muscle was measured and blockers were applied directly to them: the paralysis is lifted only if both types of \nt{GABA} receptors and the \nt{GLYC} receptors are all blocked & \cite{BrooksPeever2012} & measured \\
5 & Sleep pressure $S$ is a capacitor with two thresholds~\eqref{eq:sleeppressure}, $\tau_r=18.2$~h, $\tau_d=4.2$~h & The two-process model; the time constants are obtained from how EEG slow-wave power rises with wakefulness and falls during sleep, and the shape of the thresholds from the timing of spontaneous waking & \cite{Borbely1982,Daan1984} & theory \\
6 & The machine settles by itself into $16.1$~h awake and $8.0$~h asleep (Fig.~\ref{fig:twoprocess}) & Computed from~\eqref{eq:sleeppressure} with swinging thresholds, script \texttt{models/sleep\_models.py}: $16.1$~h awake and $7.9$--$8.0$~h asleep & --- & model \\
7 & After $40$~h without sleep $S=0.90$, but sleep lasts only $8.8$~h: the debt is repaid by depth, not by duration & Model: \texttt{models/sleep\_models.py} gives $S=0.90$ and $8.8$~h. Experiment: after $40.5$~h without sleep, in eight people EEG slow-wave power in the first recovery night is significantly above normal & \cite{Borbely1981} & model \\
8 & Physically, $S$ is adenosine & Microdialysis in cats: extracellular adenosine in one of the regions that control wakefulness rises during wakefulness, rises further with sleep deprivation, and falls during sleep. That $S$ is adenosine alone has not been shown & \cite{PorkkaHeiskanen1997,Dunwiddie2001} & hypothesis \\
9 & In slow-wave sleep the cortex swings: activity for a few hundred ms, then silence, less than once a second ($<1$~Hz, under anesthesia most often $0.3$--$0.4$~Hz) & Intracellular recording in anesthetized cats: depolarization for $0.8$--$1.5$~s and long hyperpolarization, rhythm $<1$~Hz, most often $0.3$--$0.4$~Hz; the same alternation of activity and silence is seen in natural sleep (row~1) & \cite{Steriade1993,Vyazovskiy2009} & measured \\
10 & The swings are a group with feedback and fatigue~\eqref{eq:updown}; at $b=5$ the period is $1.1$~s (a model value) and activity takes about $35\%$ of the time; at $b=1$ activity is steady (Fig.~\ref{fig:updown}) & Computation, script \texttt{models/sleep\_models.py}: period $1.11$~s, active time fraction $0.35$ (averaged over a whole number of periods); at $b=1$ the active fraction is $1.0$. The model period is shorter than measured (row~9) & --- & model \\
11 & \nt{ACET} switches off the swings and moves the cortex into steady activity & Stimulation of a nucleus of \nt{ACET} neurons in the cat blocked the slow rhythm in $79\%$ of cortical neurons and moved them into steady activity, mainly because the long hyperpolarizations disappeared. Acetylcholine suppresses the slow potassium currents that create fatigue & \cite{SteriadeAmzica1993,McCormick1992} & measured \\
12 & The $7$--$15$~Hz rhythm bursts are generated by the \nt{GLUT}--\nt{GABA} loop between the cortex and a relay node & In slices of the ferret's visual relay node, the bursts are created by rebound bursts of relay cells in response to inhibition from a neighboring \nt{GABA} nucleus, which they themselves excite & \cite{vonKrosigk1993,SteriadeMcCormickSejnowski1993} & measured \\
13 & The day's replays run fast-forward: about $20$ times faster in the hippocampus, $6$--$7$ times in the cortex & In slow-wave sleep, repeating sequences of hippocampal neurons run compressed in time. After running on a track, the sequences of place neurons were replayed in the same order in short bursts of $\sim100$~ms, compressed about $20$ times; in the cortex the compression is $6$--$7$ times & \cite{Nadasdy1999,LeeWilson2002,Euston2007} & measured \\
14 & Strengthening the alignment of replays with the cortex improves memory (a causal link) & In rats after training, stimulation locked to replays strengthened their coupling to the slow waves and rhythm bursts in the cortex: the next day memory was high, while in controls it was at chance & \cite{Maingret2016} & measured \\
15 & The purpose of replays is interleaved training of the slow memory & The theory of two learning systems and computations on artificial networks: sequential learning corrupts the old, interleaved learning does not. There is no direct experiment in the brain showing that replays are needed for exactly this & \cite{McClelland1995} & hypothesis \\
16 & After sleep synapses shrink in proportion to their size; large ones hardly change & Three-dimensional electron microscopy of $6920$ synapses in mouse cortex: the contact area after sleep is $\sim18\%$ smaller than after wakefulness, the reduction is proportional to size, and large stable synapses do not change & \cite{deVivo2017} & measured \\
17 & The main job of sleep is renormalizing the weights & The synaptic homeostasis hypothesis; rows 1 and 16 agree with it but do not prove that this is the main function & \cite{TononiCirelli2014} & hypothesis \\
18 & REM sleep is a bench test of the world model & The mode (Table~\ref{tab:sleepmodes}) is measured; that the network runs a world model is a theoretical proposal with no experiment & \cite{Hobson2009} & hypothesis \\
19 & Flushing of the brain in sleep: the question is open & One experiment: in sleep the extracellular space is $60\%$ wider and clearance is faster. Another, with a different measurement method: in sleep and under anesthesia clearance is markedly slower. The results contradict each other & \cite{Xie2013,Miao2024} & hypothesis \\
20 & Local sleep: patches of an awake animal's cortex briefly drop into silence & In rats after long wakefulness, the neurons of an individual cortical patch briefly fall silent, as in slow-wave sleep, with a slow wave in the local EEG; at such moments the rat makes more errors in a task & \cite{Vyazovskiy2011} & measured \\
\bottomrule
\end{longtable}}
